\documentclass[letterpaper,twocolumn,10pt]{article}
\usepackage{usenix2019_v3}
\usepackage{amsmath,amssymb,amsthm,enumitem,xspace,etoolbox}
\robustify{\nobreakdash}
\usepackage[capitalise,noabbrev]{cleveref}
\hypersetup{colorlinks=false}
\theoremstyle{definition}
\newtheorem{definition}{Definition}
\newtheorem{assumption}{Assumption}
\theoremstyle{plain}
\newtheorem{proposition}{Proposition}
\newcommand{\sys}{TARE\xspace}
\newcommand{\dec}[1]{\textsc{#1}}
\begin{document}
\section{Problem Model}
\label{sec:model}

The effect of a release travels along a chain: a release changes some entries, and changed entries change some
output lines. The goal is to produce every line of the new full run while running the program on as little as
possible. We make this precise.

\begin{definition}[Release]
A reference $R$ is a list of entries in file order. Each entry has a name and a \textit{strict hash} $h(e)$ of the
content that affects scoring. A new release $R'$ splits into \textit{stable} entries
$S = \{e \in R' : h(e) \in h(R)\}$ and \textit{changed} entries $\Delta = R' \setminus S$.
\end{definition}

\begin{definition}[Exact re\nobreakdash-annotation]
Let $f(Q,R)$ be the program's output on records $Q$ and reference $R$, and $\mathrm{body}(\cdot)$ the output with
comment lines removed. Given $f(Q,R)$ and $R'$, an exact re\nobreakdash-annotation returns $O$ with
$\mathrm{body}(O) = \mathrm{body}(f(Q,R'))$ byte for byte, or refuses and returns a full run $f(Q,R')$.
\end{definition}

\begin{definition}[Coupling]
A program is \textit{coupled} across entries if the line for record $q$ and entry $e$ can depend on entries other
than $e$. We distinguish four forms: a \textit{normalizer}, a global number computed from the whole reference, such as
$Z$; a \textit{side file}, an extra input file that lists entry names; an \textit{entry group}, a set of entries whose
hits compete, such as Rfam's clans; and \textit{neighbor coupling}, where the program reads bytes of the entries
next to $e$ in the file.
\end{definition}

\begin{definition}[Entry locality under a pin]
A program is \textit{entry-local under pin $\pi$} if, for every record $q$ and every subset $U \subseteq R'$ that
contains all entries reported for $q$ in the full run, running it on $(\{q\},U)$ with its normalizers set to $\pi(R')$
prints exactly $q$'s lines from the full run.
\end{definition}

Entry locality is what makes exact replay possible. It is a property of the program and its options, not something
\sys proves; \sys tests it before and during every replay and refuses when a test fails. We assume the program is
deterministic for fixed options and thread count, and that every input it reads is named on its command line.
\sys treats the program as a black box at run time and never reads its source code. We, the authors, did read the
source code of HMMER, Infernal, and DIAMOND, to design experiments and to explain their outcomes; none of that knowledge
is built into \sys.

The gate applies three rules in order and stops at the first that fires:
\begin{equation}
D = \left\{\begin{array}{@{}lll@{}}
\dec{refuse} & \text{no contract} & \text{(R1)}\\
\dec{refuse\_pin} & \text{size-dependent, no pin} & \text{(R2)}\\
\dec{refuse\_pin} & \text{in-situ check fails} & \text{(R3)}\\
\dec{authorize} & \text{otherwise} &
\end{array}\right.
\label{eq:gate}
\end{equation}
R1 fires when the output does not split or a column fits no hypothesis. R2 fires when a column depends on the database
size and no pin is available; \sys never guesses. R3 replays the probe sample against its own hits with the pin set and
requires every byte to match the whole run. Every refusal runs the program in full; \dec{refuse\_pin} differs from
\dec{refuse} only in what it reports: the output does decompose, so declaring a working pin could enable replay. Before replaying, the gate also falls back
to a full run when $f_\Delta$, the changed fraction of the reference, exceeds the break\nobreakdash-even point $(1-c)/a$ of the
cost model (the cost model); with no timing history it uses 0.80.
depends only on that protein's own hits, including the column widths and the per\nobreakdash-protein statistics. As a tripwire, if
a replay ever reports a protein--family pair outside the candidates, \sys refuses (R4), because a full run could
contain pairs the cache never saw.
\subsection{Stage 6: Check or Refuse}
\label{sec:check}

The replay assumes the program is entry\nobreakdash-local. The check tries to prove that assumption wrong on the actual run. It
reruns the replay with \textit{pads}, families that were deliberately left out, added back to the subset. If the
program is entry\nobreakdash-local, pads change nothing: no candidate line changes and no pad family is reported. If anything
changes, the replay depended on what it left out, and \sys refuses (R5).

\sys uses two kinds of pads. The \textit{neighbor pad} adds, for every replayed protein, the families that sit next to
its candidates in the database file; it catches neighbor coupling by construction, the failure we found in DIAMOND
(the evaluation). \textit{Random pads} add three random samples of left\nobreakdash-out families, each as large as the
candidate set, for a fixed sample of $s=50$ proteins.
% source: MEASURED pipeline/docs/EXACT_REPLAY_PROTOCOL.md; TEXT sample_s=50
They test for dependencies of unknown form. If an unknown
dependency affects $k$ of $n$ proteins and a random pad includes the responsible family with probability $q$, the check
misses it with probability approximately
\begin{equation}
P_{\text{miss}} \approx h + (1-h)(1-q)^3, \qquad h = \binom{n-k}{s}\Big/\binom{n}{s},
\label{eq:miss}
\end{equation}
% source: MEASURED pipeline/results/reference_diamond/dx3.json; LEN pad_seeds=3; NOMINAL one=1
where $h$ is the chance that the protein sample contains no affected protein. The approximation treats each pad as a
single detection event; when several affected proteins fall in the sample, each pad has more chances to detect, so the
true miss probability is lower. The neighbor pad needs no such formula: it contains every candidate's file\nobreakdash-order
neighbor by construction. The check is a test, not a proof. As a final safeguard, an audit reruns the stock program on a
random sample of $s$ proteins and refuses if any line differs (R6). By default
$s = \lceil \ln\varepsilon / \ln(1-\kappa) \rceil = 299$ for $\varepsilon = 0.05$ and $\kappa = 0.01$, capped at the
record count: if at least a fraction $\kappa$ of the records is affected, the audit misses all of them with probability
at most $\varepsilon$. This is the classical zero\nobreakdash-defect audit sample, not a new bound. The measured runs in
the evaluation that predate it used an earlier audit, each record included with probability 0.02 and at least 20 in
total. An optional first\nobreakdash-use mode instead compares the first two updates of a new program with a full run.
\end{document}
